%==========================
%          ABSTRACT
%==========================
\chapter*{\hfill{\centering Abstract}\hfill}
\addcontentsline{toc}{chapter}{Abstract}  %Retirar comentario si quiere que aparesca en la lista del Indice general
\vspace{0.5cm}
\justify
\normalsize

This research work presents an analytical and perturbative analysis of photon geodesic dynamics and shadow structure in the Kerr-Newman spacetime under the weak electric charge approximation ($\delta \equiv \frac{Q^2}{4\pi\varepsilon_0 G M^2} \ll 1$), rigorously formulated in International System (SI) units and employing the Minkowski metric signature convention $(+,-,-,-)$. The effective radial potential $\mathcal{R}(r)$ is derived step by step using the Hamilton-Jacobi formalism and the action function $S$, yielding the conserved constants of motion: energy $E$, azimuthal angular momentum $L_z$, and Carter's constant $\mathcal{K}$. Defining the reduced impact parameters $\Phi \equiv \frac{c L_z}{E}$ and $\mathcal{Q} \equiv \frac{c^2 \mathcal{K}}{E^2}$, and imposing the unstable spherical photon orbit conditions $\mathcal{R}(r_{\mathrm{ph}}) = 0$ and $\left.\frac{d\mathcal{R}}{dr}\right|_{r_{\mathrm{ph}}} = 0$, Kobialko's methodology is extended to non-diagonal rotational metrics via the critical parameter $b_c^2 = \mathcal{Q}_c + \Phi_c^2$. We derive the perturbative expansion of the equivalent effective radius $R_{\mathrm{eff}}(\delta)/R_{\mathrm{eff}}(0) = 1 - \frac{1}{6}\delta - \frac{7}{216}\delta^2$, the second-order coefficient $D = 1/27$ in $R_{\mathrm{eff}}^2$ isolating the unique electrovacuum charge signature from modified gravity models, and the displacement ratio of the first visibility null $u_1(\delta)/u_1(0) = 1 + \frac{1}{6}\delta + \frac{13}{216}\delta^2$. Furthermore, it is shown that while in the static Schwarzschild limit Sneppen's thinning factor ($f_{\mathrm{thin}} = e^{-2\pi} \approx 1/535$) renders the $n=1$ subring unobservable, in Kerr-Newman the angular momentum degrades the Lyapunov exponent ($\gamma_\phi < 1$), orbitally widening the first photon subring. Finally, null geodesics are integrated using the fourth-order Runge-Kutta (RK4) algorithm in Python, and the theoretical interferometric visibility profiles $|V(u)|$ in the Fourier domain are contrasted with Event Horizon Telescope (EHT) observations for $\text{Sgr A}^*$ (2024-D02-01 data release) and $\text{M87}^*$ (2024-D01-01 data release). For $\text{Sgr A}^*$, the Kerr-Newman model with $\delta \approx 0.290$ shifts the first visibility null by $+5.42\%$, improving the goodness-of-fit by $\Delta\chi^2_\nu = -4.32$ ($\Delta\chi^2 = -40\,166$) compared to Schwarzschild; whereas for $\text{M87}^*$, the 1D symmetric model degrades ($\Delta\chi^2_\nu = +10.74$) due to contamination from extended relativistic jet emission. It is concluded that the fit in $\text{Sgr A}^*$ acts as a geometric scale compensator within the 1D symmetric model rather than a direct physical charge detection, given the dominance of current systematic uncertainties in the mass-to-distance ratio ($\sim 11.8\%$).


%=================================
% Keywords
%=================================
\vspace{0.5cm}
\raggedright 
{\bf \textit{Keywords}:} Kerr-Newman black hole, Weak electric charge, Photon sphere, Effective shadow radius, Carter's constant, VLBI Interferometry, Event Horizon Telescope, Lyapunov exponent.

